\documentclass[11pt,a4paper]{article}

\usepackage[T1]{fontenc}
\usepackage[utf8]{inputenc}
\usepackage[czech]{babel}
\usepackage{lmodern}
\usepackage{amsmath,amssymb,amsthm,mathtools}
\usepackage{microtype}
\usepackage{geometry}
\usepackage[hidelinks]{hyperref}

\geometry{margin=2.5cm}

\newtheorem{theorem}{Věta}
\newtheorem{lemma}[theorem]{Lemma}
\newtheorem{proposition}[theorem]{Tvrzení}
\newtheorem{corollary}[theorem]{Důsledek}
\theoremstyle{definition}
\newtheorem{definition}[theorem]{Definice}
\theoremstyle{remark}
\newtheorem{remark}[theorem]{Poznámka}
\theoremstyle{plain}
\newtheorem{conjecture}[theorem]{Otevřený problém}

\newcommand{\E}{\mathbb E}
\newcommand{\Prb}{\mathbb P}
\newcommand{\cK}{\mathcal K}
\newcommand{\Law}{\operatorname{Law}}
\newcommand{\DOWN}{\operatorname{DOWN}}
\newcommand{\UP}{\operatorname{UP}}
\newcommand{\st}{\mathrm{st}}

\title{Quartic Zero-Repair:\\bezpodmínečný sandwich a otevřený logaritmus}
\author{}
\date{16.~září 2026}

\begin{document}
\maketitle

\begin{abstract}
Dokazujeme bezpodmínečný sandwich pro kvartický zero-repair problém
$(P,Q)$ na abecedě $\{0,1,2,3,4\}$:
\[
\kappa_0\log n-o(\log n)
\le\mathcal B_n
\le\frac{473}{56}\,n^{1/4}.
\]
Horní mez je explicitní konstrukce s negativně binomickou referencí.
Dále dokazujeme Walshovu identitu stupně $4$, causal thinning na parity
core $(E,O)$ a blokový most
$0\le\mathcal B_{2n}(E,O)-2\mathcal B_n(P_{\mathrm b},Q_{\mathrm b})\le 4$.
Z toho plyne, že $\mathcal B_n=O(\log n)$ je ekvivalentní témuž odhadu
pro coarse pair $(P_{\mathrm b},Q_{\mathrm b})$ až na konstantu.
\textbf{Tvrzení $\mathcal B_n=O(\log n)$ v tomto článku nedokazujeme.}
Poslední chybějící krok je jeden prefix-kompatibilní, phase-indexed fan
realizující dyadický boundary-flow
$\mathcal L_n(\delta)\to\mathcal L_{2n}(\frac12\delta+C\delta^2)$.
\end{abstract}

\tableofcontents
\newpage

\section{Úvod}

Exactní zero-repair strom je kauzální protokol, který v každém kroku
losuje čerstvý vstup podle $Q$ a veřejný výstup podle $P$, přičemž
skrytá rezerva $J$ zůstává nezáporná. Cena hloubky $n$ je infimum
počáteční střední rezervy. Dolní mez $\kappa_0\log n$ je známa.
Otázka, zda $\mathcal B_n=\Theta(\log n)$, zůstává otevřená.

Tento text je článek, ne pracovní ledger. Obsahuje pouze věty s důkazem
a na konci explicitně označený otevřený problém. Žádné numerické
extrapolace se zde neprohlašují za věty.

Hlavní výsledek je sandwich. Cesta k logaritmu nahoře vede přes thinning
a blokový most na coarse pair, a končí u dyadického rozbalení boundary
packetu mezi škálami $n\mapsto 2n$. Faktor $1/2$ je survival přes slupku,
nikoli větvení pěti dětí. Prefixový zámek zakazuje různé donory na
společném prefixu. Ten fan není sestrojen.

\section{Model}

Nechť $B_h$ jsou pravděpodobnostní zákony na $\mathbb N_0$ indexované
veřejnými historiemi $h\in\{0,1,2,3,4\}^{\le n}$. Exactní zero-repair
strom splňuje v každém vnitřním uzlu
\[
Q*B_h=\sum_{x=0}^4 P_x\,\delta_x*B_{hx}.
\tag{ZR}
\]
Množinu kořenových zákonů hloubky $n$ označme $\cK_n$ a položme
\[
\mathcal B_n:=\inf_{B\in\cK_n}\E_B J.
\]

Pro generující funkce platí přesná kvartická identita
\[
G_P(z)-G_Q(z)=\frac{(1-z)^4}{64}.
\tag{1}
\]

\begin{theorem}[dříve dokázaná dolní mez]
Platí
\[
\liminf_{n\to\infty}\frac{\mathcal B_n}{\log n}\ge \kappa_0,
\qquad
\kappa_0=
\frac{\log(594184/543559)}
{\log 16\,\log(524288/74273)}
\approx0.0164347253.
\]
\end{theorem}

\section{Explicitní horní mez $O(n^{1/4})$}

Položme $D=\delta_0-\delta_1$. Z~(1) plyne
\[
P-Q=\frac1{64}D^{*4}.
\tag{2}
\]
Pro $a>0$ definujme geometrický zákon
\[
\gamma_a(j)=\frac1{1+a}\left(\frac a{1+a}\right)^j,
\qquad j\ge0,
\]
a
\[
C_a:=\gamma_a^{*4}.
\]
Dále položme
\[
\varepsilon:=\frac{25}{7a^3},
\qquad
E_\varepsilon:=(1-\varepsilon)\delta_0+\varepsilon\delta_1,
\qquad
R_r:=C_a*E_\varepsilon^{*r}.
\tag{3}
\]
Je-li $a^3\ge25/7$, pak $0\le\varepsilon\le1$ a
\[
\E R_r=4a+r\varepsilon.
\tag{4}
\]

\begin{lemma}[koeficientová nerovnost]
Pro každé $a>0$ platí koeficient po koeficientu
\[
D^{*3}*C_a+\frac{25}{a^3}C_a\ge0.
\tag{5}
\]
Tedy pro každé $r\ge0$
\[
D^{*3}*R_r\ge-\frac{25}{a^3}R_r.
\tag{6}
\]
\end{lemma}

\begin{proof}
Z geometrické řady
\[
D*\gamma_a=\frac1a(\delta_0-\gamma_a).
\]
Proto
\[
D^{*3}*C_a
=
\frac1{a^3}
\left(
\gamma_a-3\gamma_a^{*2}+3\gamma_a^{*3}-\gamma_a^{*4}
\right).
\]
Po přičtení $25C_a/a^3$ stačí ověřit nezápornost
\[
\gamma_a-3\gamma_a^{*2}+3\gamma_a^{*3}+24\gamma_a^{*4}.
\]
Pišme $\delta=(1+a)^{-1}$. V koeficientu $j\ge0$, po vydělení
kladným $\gamma_a(j)$, dostáváme
\[
1-3\delta(j+1)
+\frac32\delta^2(j+1)(j+2)
+4\delta^3(j+1)(j+2)(j+3).
\]
Pro $u=\delta(j+1)\ge0$ je tento výraz zdola omezen
\[
1-3u+\frac32u^2+4u^3
=
(u+1)(2u-1)^2+\frac32u^2\ge0.
\]
To dokazuje (5). Konvoluce s nezápornou mírou zachovává
koeficientovou nezápornost, a tedy (6).
\end{proof}

\begin{lemma}[dominance referencí]
Pro $r\ge1$ platí
\[
Q*R_r\succeq_{\st}P*R_{r-1}.
\tag{7}
\]
\end{lemma}

\begin{proof}
Z $R_r=R_{r-1}*E_\varepsilon$ a (2)
\[
Q*R_r-P*R_{r-1}
=
-\frac1{64}D^{*4}*R_{r-1}
-\varepsilon D*Q*R_{r-1}.
\]
Částečný součet koeficientů $D*M$ do $k$ je $M(k)$, proto
\[
F_{Q*R_r}(k)-F_{P*R_{r-1}}(k)
=
-\frac1{64}(D^{*3}*R_{r-1})(k)
-\varepsilon(Q*R_{r-1})(k).
\]
Protože $Q_0=7/64$,
\[
Q*R_{r-1}\ge\frac7{64}R_{r-1}
\]
koeficient po koeficientu. Použitím (6) je pravá strana nejvýše
\[
\left(
\frac{25}{64a^3}-\frac{7\varepsilon}{64}
\right)R_{r-1}(k)=0.
\]
Tedy $F_{Q*R_r}\le F_{P*R_{r-1}}$.
\end{proof}

\begin{theorem}[explicitní all-horizon konstrukce]
Pro každé $n\ge1$ a každé $a>0$ s $a^3\ge25/7$ platí
\[
\boxed{
\mathcal B_n\le 4a+\frac{25n}{7a^3}.
}
\tag{8}
\]
Zvláště pro $a=2n^{1/4}$,
\[
\boxed{
\mathcal B_n\le\frac{473}{56}\,n^{1/4}.
}
\tag{9}
\]
\end{theorem}

\begin{proof}
Začněme z $B_\varnothing=R_n$. Budeme induktivně udržovat
\[
B_h\succeq_{\st}R_r,
\]
kde $r$ je počet zbývajících kroků. Je-li tato podmínka splněna,
pak monotonicita konvoluce a (7) dávají
\[
\mu:=Q*B_h\succeq_{\st}\alpha:=P*R_{r-1}.
\]
Existuje tedy monotónní vazba $(U,S)$ s
\[
U\sim\alpha,\qquad S\sim\mu,\qquad U\le S\quad\text{a.s.}
\]
Rozložme $U=X+Z$ podle produktu
\[
\Prb(X=x,Z=z)=P_xR_{r-1}(z).
\]
Použijeme podmíněný zákon $X$ daný $U$ a~$S$ a definujeme
\[
J':=S-X.
\]
Protože $X\le U\le S$, je $J'\ge0$. Veřejný marginál $X$ je přesně
$P$. Dále
\[
J'=S-X\ge U-X=Z.
\]
Podmíněně na každém $X=x$ má $Z$ zákon $R_{r-1}$, takže každý
posterior $B_{hx}$ splňuje
\[
B_{hx}\succeq_{\st}R_{r-1}.
\]
Indukce tedy pokračuje až do hloubky $n$.

Počáteční účet je podle (4)
\[
\E B_\varnothing
=
4a+n\varepsilon
=
4a+\frac{25n}{7a^3},
\]
což dává (8). Pro $a=2n^{1/4}$ dostáváme
\[
8n^{1/4}+\frac{25}{56}n^{1/4}
=
\frac{473}{56}n^{1/4}.
\]
\end{proof}

\section{Předplacené doplnění}

Pro slovo $u=(u_1,\dots,u_r)$ pišme
\[
p(u):=\prod_{i=1}^r P_{u_i}.
\]

\begin{theorem}[předplacené doplnění]
Nechť je dán exactní strom hloubky $n$ s kořenem $B$. Předpokládejme,
že z každého jeho listu $h$ existuje pomocný kauzální $r$-krokový
protokol s čerstvým $Q$, nezápornou rezervou a veřejným zákonem slov
$R_h$, přičemž pro společné $M\ge1$
\[
R_h(u)\le M p(u)
\qquad\text{pro všechna }h,u.
\tag{10}
\]
Pak existuje exactní strom hloubky $n+r$ s kořenem
\[
B*\left(\frac1M\delta_0+
\left(1-\frac1M\right)\delta_{4r}\right),
\]
a počáteční mean se zvětší právě o
\[
\boxed{
4r\left(1-\frac1M\right).
}
\tag{11}
\]
\end{theorem}

\begin{proof}
Položme $\eta=1/M$ a $\epsilon=1-\eta$. Pro každý list definujme
\[
D_h(u):=
\frac{p(u)-\eta R_h(u)}{\epsilon}
\]
pro $\epsilon>0$; při $M=1$ není opravná větev potřeba.
Z (10) plyne $D_h(u)\ge0$ a
\[
\sum_uD_h(u)=1,\qquad
\eta R_h+\epsilon D_h=P^{\otimes r}.
\]
Na začátku nezávisle losujeme Bernoulliho značku $I$ s
$\Prb(I=1)=\epsilon$ a přidáme rezervu $4rI$.
Prvních $n$ kroků proběhne původní exactní strom.

Je-li $I=0$, použije se pomocný protokol $R_h$.
Je-li $I=1$, použije se veřejný zákon $D_h$; předplacená rezerva
$4r$ stačí na libovolné slovo délky $r$, protože každý výstup je
nejvýše $4$. V obou větvích se vstupy losují čerstvě podle $Q$.
Po zprůměrování přes $I$ je veřejný zákon následujícího slova
přesně $P^{\otimes r}$. Cena značky je
\[
\E(4rI)=4r\epsilon=4r(1-1/M).
\]
\end{proof}

\section{Explicitní Q-fresh vazba}

Pro zákon $B$ pišme
\[
p=B(0),\qquad q=B(1).
\]

\begin{lemma}[Q-fresh vazba]
Je-li $p\le1/4$, položme
\[
\theta=-\frac{3p}{1-p}\in[-1,0].
\]
Na stavu $J=0$ použijme coupling
\[
C^{(0)}=\frac1{64}
\begin{pmatrix}
7&4&0&0&0\\
0&16&0&0&0\\
0&0&10&4&2\\
0&0&0&16&0\\
0&0&0&0&5
\end{pmatrix},
\]
a na $J\ge1$
\[
C^{(\theta)}=\frac1{64}
\begin{pmatrix}
7&1+\theta&0&0&0\\
0&16&0&0&0\\
0&3-\theta&10&3+\theta&0\\
0&0&0&16&0\\
0&0&0&1-\theta&7
\end{pmatrix}.
\]
Oba couplingy mají sloupce $Q$, první má podporu $X\le Y$ a druhý
$X-Y\le1$. Po zprůměrování přes skrytý stav má $X$ přesně zákon $P$.
\end{lemma}

\begin{proof}
Řádkové zákony jsou
\[
R^{(0)}=P+\frac3{64}(\delta_0-\delta_4),
\qquad
R^{(\theta)}=P+\frac{\theta}{64}(\delta_0-\delta_4).
\]
Směs podle pravděpodobností $p$ a $1-p$ je $P$ právě když
\[
3p+(1-p)\theta=0.
\]
Podmínky podpory dávají nezápornost nové rezervy.
\end{proof}

Položme
\[
a=\frac{1-4p}{1-p},\qquad
c=\frac{1+2p}{1-p},\qquad
u=\frac3{1-p},\qquad
v=\frac{3(1-2p)}{1-p}.
\tag{12}
\]
Na kladné rezervě mají aktivní děti tvar
\[
T_0=\frac78I+\frac a8\UP,
\qquad
T_2=\frac u{16}\DOWN+\frac{10}{16}I+\frac v{16}\UP,
\qquad
T_4=\frac c8\DOWN+\frac78I,
\tag{13}
\]
a $T_1=T_3=I$. Platí $a\ge0$ právě pro $p\le1/4$.

Jednokrokové nulové hmotnosti jsou
\[
p'_0=\frac78p,\qquad
p'_2=\frac{10p+uq}{16},\qquad
p'_4=\frac{5p+cq}{8}.
\tag{14}
\]

\begin{lemma}[jednokroková face]
Pro $p<1$ položme $\eta=q/(1-p)$. Podmínka $p'_2\le1/4$ je ekvivalentní
\[
\eta\le\frac{4-10p}{3}.
\]
Podmínka $p'_4\le1/4$ je
\[
\eta\le\frac{2-5p}{1+2p},
\]
a na $0\le p\le1/4$ je druhá mez slabší.
\end{lemma}

Definujme tedy
\[
\cK:=
\left\{
0\le p\le\frac14,\quad
0\le\eta\le
\min\left(1,\frac{4-10p}{3}\right)
\right\}.
\tag{15}
\]

\begin{lemma}
Platí
\[
T_0(\cK)\subset\cK.
\tag{16}
\]
\end{lemma}

\begin{proof}
Pro $T_0$
\[
p'=\frac78p,\qquad
\eta'=
\frac{4p+7\eta(1-p)}{8-7p}.
\]
Mapování je rostoucí v $\eta$. Na stěně
$\eta=(4-10p)/3$ se požadovaná nerovnost redukuje na
\[
35p^2+48p-16\le0.
\]
Kladný kořen je
\[
\frac{4(-6+\sqrt{71})}{35}>\frac14,
\]
takže nerovnost platí na $[0,1/4]$.
Na víčku $\eta=1$ je
\[
\eta'=\frac{7-3p}{8-7p}\le1
\]
pro $p\le1/4$.
\end{proof}

\section{Waringova a Beta reprezentace}

\begin{definition}
Pro $k\ge0$ nechť
\[
\Prb(V_k\ge j)=\frac{k+1}{k+j},
\qquad j\ge1,
\]
a $W_{r,k}:=\Law(r+V_k)$.
\end{definition}

\begin{lemma}[Waringova identita]
\[
\boxed{
\Law(V_k-1)
=
\frac1{k+2}\delta_0+
\frac{k+1}{k+2}V_{k+1}.
}
\tag{17}
\]
\end{lemma}

\begin{proof}
Pro $j\ge1$
\[
\Prb(V_k-1\ge j)
=
\frac{k+1}{k+j+1}
=
\frac{k+1}{k+2}\frac{k+2}{k+j+1}.
\]
\end{proof}

Pro $0\le t<1$ definujme geometrický zákon na $\{1,2,\dots\}$
\[
g_t(j)=(1-t)t^{j-1}
\]
a $E_{r,t}:=\delta_r*g_t$.

\begin{lemma}[Beta-mix]
\[
\boxed{
W_{r,k}=(k+1)\int_0^1 t^kE_{r,t}\,dt.
}
\tag{18}
\]
\end{lemma}

\begin{proof}
Pro $j\ge1$
\[
\Prb(E_{r,t}\ge r+j)=t^{j-1},
\]
a tedy
\[
(k+1)\int_0^1 t^{k+j-1}\,dt
=
\frac{k+1}{k+j}.
\]
\end{proof}

\begin{lemma}[paměť na hraně]
\[
\boxed{
\DOWN E_{0,t}=(1-t)\delta_0+tE_{0,t}.
}
\tag{19}
\]
Pro $r\ge1$
\[
\DOWN E_{r,t}=E_{r-1,t},
\qquad
\UP E_{r,t}=E_{r+1,t}.
\]
\end{lemma}

Definujme konkrétní Waringův zákon
\[
B^{(3)}
=
\frac14\delta_0+\frac34W_{0,3}
=
\frac14\delta_0+\int_0^1 E_{0,t}\,3t^3\,dt.
\tag{20}
\]
Jeho prvních několik hmotností je
\[
B^{(3)}(0)=\frac14,\qquad
B^{(3)}(1)=\frac3{20},\qquad
B^{(3)}(2)=\frac1{10},
\]
a pro $j\ge1$
\[
B^{(3)}(j)=\frac{3}{(j+3)(j+4)},
\qquad
\Prb(J\ge j)=\frac3{j+3}.
\tag{21}
\]

\begin{lemma}[kladný shape cone]
Uvažujme rozklad
\[
B=p\delta_0+\sum_{r\ge0}\int_0^1E_{r,t}\,\nu_r(dt).
\]
Je-li
\[
\nu_r=e_r\delta_0+t^3f_r(t)\,dt,
\qquad f_r(t)\in\mathbb R_{\ge0}[t],
\tag{22}
\]
pak operátory (13), dokud $p\le1/4$, tento tvar zachovávají.
\end{lemma}

\begin{proof}
Všechny koeficienty v (13) jsou nezáporné. Operace $\UP$ a vnitřní
$\DOWN$ mění pouze index $r$. Hraniční $\DOWN$ podle (19) násobí
kontinuální část faktorem $t$ a přidává nezáporný atom v $t=0$.
Tedy $t^3f_r(t)$ zůstává tvaru $t^3\widetilde f_r(t)$ s
$\widetilde f_r\ge0$ koeficient po koeficientu.
\end{proof}

Definujme
\[
h_r:=\int_0^1(1-t)\,\nu_r(dt).
\tag{23}
\]
Pak
\[
q=B(1)=h_0.
\]

\begin{lemma}[sandwich pro druhou hmotnost]
Na kuželu (22)
\[
\boxed{
h_1+\frac23(q-e_0)
\le B(2)
\le h_1+(q-e_0).
}
\tag{24}
\]
Navíc
\[
e_r\le h_r.
\tag{25}
\]
\end{lemma}

\begin{proof}
V ekvivalentních Waringových souřadnicích
\[
\nu_r=e_r\delta_0+\sum_{k\ge3}\omega_{r,k}(k+1)t^k\,dt,
\qquad \omega_{r,k}\ge0.
\]
Pak
\[
q-e_0=\sum_{k\ge3}\frac{\omega_{0,k}}{k+2},
\]
zatímco
\[
B(2)-h_1
=
\sum_{k\ge3}
\frac{k+1}{(k+2)(k+3)}\omega_{0,k}.
\]
Pro $k\ge3$
\[
\frac23\le\frac{k+1}{k+3}<1,
\]
což dává (24). Dále
\[
h_r=e_r+\sum_{k\ge3}\frac{\omega_{r,k}}{k+2}\ge e_r.
\]
\end{proof}

\begin{lemma}[fronty pod $T_0$]
Souřadnice $(p,q,h_1,e_0)$ jsou pod $T_0$ uzavřené:
\[
p'=\frac78p,\qquad
q'=\frac p2+\frac78q,
\]
\[
h_1'=\frac78h_1+\frac a8q,
\qquad
e_0'=\frac p2+\frac78e_0.
\tag{26}
\]
Residuum $\sigma:=q-e_0$ splňuje
\[
\boxed{\sigma'=\frac78\sigma.}
\tag{27}
\]
\end{lemma}

Pro orbitu $T_0^nB^{(3)}$ tedy
\[
p_n=\frac14\left(\frac78\right)^n,
\qquad
q_n=
\left(\frac78\right)^n
\left(\frac3{20}+\frac n7\right).
\tag{28}
\]

\section{Exactní děti zákona $B^{(3)}$}

\begin{lemma}
Pro vazbu výše platí
\[
T_0B^{(3)}
=
\frac7{32}\delta_0+
\frac18\delta_1+
\frac{21}{32}W_{0,3},
\tag{29}
\]
\[
T_4B^{(3)}
=
\frac{31}{160}\delta_0+
\frac{21}{32}W_{0,3}
+\frac3{20}W_{0,4},
\tag{30}
\]
a
\[
T_2B^{(3)}
=
\frac{31}{160}\delta_0+
\frac1{16}\delta_1+
\frac1{32}\delta_2+
\frac{15}{32}W_{0,3}
+\frac3{20}W_{0,4}
+\frac3{32}W_{1,3}.
\tag{31}
\]
\end{lemma}

\begin{proof}
V $B^{(3)}$ je $p=1/4$, a tedy
\[
a=0,\qquad c=2,\qquad u=4,\qquad v=2.
\]
Dosazení do (13) a použití (17) dává (29)--(31).
\end{proof}

\begin{corollary}
Platí
\[
T_4B^{(3)}\succeq_{\st}B^{(3)}.
\tag{32}
\]
\end{corollary}

\begin{proof}
Pro $j\ge1$
\[
S_{T_4B^{(3)}}(j)
=
\frac{3(9j+34)}{8(j+3)(j+4)},
\]
zatímco $S_{B^{(3)}}(j)=3/(j+3)$. Proto
\[
\frac{S_{T_4B^{(3)}}(j)}{S_{B^{(3)}}(j)}
=
\frac{9j+34}{8j+32}\ge1.
\]
\end{proof}

Pro $T_2$ a $j\ge3$ platí přesně
\[
\frac{S_{T_2B^{(3)}}(j)}{S_{B^{(3)}}(j)}
=
\frac{8j^2+47j+64}{8j^2+48j+64}
=
1-\frac{j}{8(j+2)(j+4)}.
\tag{33}
\]

\section{Minimaxový one-fan}

Nechť $B$ je libovolný pravděpodobnostní zákon a
\[
\mu:=Q*B.
\]
Předpokládejme, že one-fan existuje, tj.
\[
P\preceq_{\st}\mu
\quad\Longleftrightarrow\quad
F_\mu(k)\le F_P(k)\quad(k\ge0).
\tag{34}
\]
One-fan je coupling $\pi(x,s)$ s marginály $P$ a $\mu$ a podporou
$x\le s$. Příslušné dítě splňuje
\[
B_x(j)=\frac{\pi(x,x+j)}{P_x},
\]
takže
\[
B_x(0)=\frac{\pi(x,x)}{P_x}.
\]

\begin{definition}
Definujme
\[
\tau(B):=
\inf_\pi\max_{0\le x\le4}B_x(0),
\tag{35}
\]
kde infimum běží přes všechny one-fany.
\end{definition}

\begin{theorem}[přesný minimax diagonály]
Platí
\[
\boxed{
\tau(B)=
\max_{0\le k\le4}
\left(
\frac{F_\mu(k)-F_P(k-1)}{P_k}
\right)_+,
}
\tag{36}
\]
kde $F_P(-1):=0$.
Navíc jediný kvantilový coupling současně minimalizuje každou
diagonální hmotnost $\pi(k,k)$.
\end{theorem}

\begin{proof}
Fixujme $k$. Hmota sloupců $s\le k$ je $F_\mu(k)$. Řádky
$0,\dots,k-1$ mohou absorbovat nejvýše $F_P(k-1)$. Přebytek
\[
\bigl(F_\mu(k)-F_P(k-1)\bigr)_+
\]
proto musí skončit v řádku $k$. Ze sloupců $s\le k$ je pro řádek
$k$ povolen pouze bod $(k,k)$. Tedy pro každý fan
\[
\pi(k,k)\ge
\bigl(F_\mu(k)-F_P(k-1)\bigr)_+.
\tag{37}
\]

Nyní vezměme kvantilový coupling. Na jednotkovém intervalu mají
události $X=k$ a $S=k$ intervaly
\[
(F_P(k-1),F_P(k)],
\qquad
(F_\mu(k-1),F_\mu(k)].
\]
Z (34)
\[
F_\mu(k-1)\le F_P(k-1),
\qquad
F_\mu(k)\le F_P(k).
\]
Délka jejich průniku je tedy přesně
\[
\pi(k,k)
=
\bigl(F_\mu(k)-F_P(k-1)\bigr)_+.
\tag{38}
\]
Kvantilový coupling má současně podporu $X\le S$. Dolní meze (37)
jsou proto dosaženy pro všechna $k$ zároveň. Po vydělení $P_k$
dostáváme (36).
\end{proof}

\section{Vyvážený kužel}

Pišme
\[
p=B(0),\quad q=B(1),\quad r=B(2),\quad s=B(3),\quad t=B(4).
\]

\begin{theorem}[jednokroková kontrakce nulové hmotnosti]
Následující čtyři nerovnosti
\[
13p+7q\le8,
\tag{39a}
\]
\[
23p+27q+7r\le24,
\tag{39b}
\]
\[
43p+37q+27r+7s\le40,
\tag{39c}
\]
\[
57p+57q+37r+27s+7t\le56
\tag{39d}
\]
jsou ekvivalentní podmínce
\[
\tau(B)\le\frac78p.
\tag{40}
\]
Je-li (39a)--(39d) splněno, existuje exactní one-fan, pro který
\[
\boxed{
B_x(0)\le\frac78B(0)
\qquad\text{pro všechna }x=0,\dots,4.
}
\tag{41}
\]
Pro $x=0$ je rovnost nevyhnutelná.
\end{theorem}

\begin{proof}
Pro $k=0$ je
\[
\frac{F_\mu(0)}{P_0}
=
\frac{Q_0p}{P_0}
=
\frac78p.
\]
Pro $k=1,2,3,4$ dosadíme $\mu=Q*B$ do (36) a požadujeme, aby
každý člen byl nejvýše $7p/8$. Po vynásobení $64$ dostáváme
právě (39a)--(39d). Tyto nerovnosti zároveň dávají
$F_\mu(k)\le F_P(k)$, takže one-fan existuje. Kvantilový fan z
předchozí věty realizuje (41).
\end{proof}

\begin{corollary}[zákon $B^{(3)}$ splňuje vyvážené nerovnosti]
Pro $B^{(3)}$
\[
p=\frac14,\quad
q=\frac3{20},\quad
r=\frac1{10},\quad
s=\frac1{14},\quad
t=\frac3{56}.
\]
Levé strany (39a)--(39d) jsou po řadě
\[
\frac{43}{10},\qquad
\frac{21}{2},\qquad
\frac{39}{2},\qquad
\frac{1613}{56},
\]
a tedy jsou ostře menší než
\[
8,\qquad24,\qquad40,\qquad56.
\]
Proto existuje one-fan z $B^{(3)}$ se
\[
B_x(0)\le\frac7{32}
\qquad(x=0,\dots,4).
\tag{42}
\]
Přesná minimaxová hodnota je
\[
\boxed{\tau(B^{(3)})=\frac7{32}.}
\tag{43}
\]
\end{corollary}

\section{Bezpodmínečný sandwich}

Z tvrzení výše plyne bez dalších hypotéz
\[
\boxed{
\kappa_0\log n-o(\log n)
\le
\mathcal B_n
\le
\frac{473}{56}\,n^{1/4}.
}
\tag{44}
\]
Rovnice (17)--(33) jsou přesné identity pro uvedenou Q-fresh vazbu.
Věty (36) a (39)--(43) jsou přesné jednokrokové optimalizační výroky.
Žádné tvrzení v tomto dokumentu neprohlašuje
$\mathcal B_n=O(\log n)$.

\section{$T_0$ ve čtyřech souřadnicích a $\mu=1$}

Residuum $\sigma:=q-e_0$ splňuje $\sigma'=7\sigma/8$. Čtveřice
$(p,q,h_1,e_0)$ je pod $T_0$ uzavřená.

\begin{theorem}[$T_0$ drží $q+h_1\le q_{\max,2}$]
Nechť $q_{\max,2}(p)=(1-p)\min\bigl(1,(4-10p)/3\bigr)$.
Je-li $q+h_1\le q_{\max,2}(p)$ a $p\in[0,1/4]$, totéž platí po $T_0$.
Nejhorší případ je $q=V$, $h_1=0$. Slack proti $q_{\max,2}(7p/8)$ je
\[
\frac{p(56-195p)}{96}\ge0
\qquad\text{na }[0,1/4],
\]
protože $195/4<56$. V $p=1/4$ je slack $29/1536$.
\end{theorem}

\begin{remark}
$\mu>1$ v nejhorším případě neplatí už v $p=0$. $T_2$ se $\mu=1$
neuzavírá: na stěně $\mathcal K$ s $(p,q)=(1/10,9/10)$ a $B(2)=0$ je
$q'/q_{\max,2}(p')=47/30>1$. Separátní $\{h\text{ klesá}\}$ a
$\{h_1\le q\}$ padají na $T_0^{18}(B^{(3)})$.
\end{remark}

Pro děti $B^{(3)}$ platí $B(2)$
\[
T_0:\ \tfrac7{80},\qquad
T_4:\ \tfrac{59}{560},\qquad
T_2:\ \tfrac{73}{560}.
\]

\section{Geometrie předplacení}

Při hypotéze $M_n\le 1+a/n$ na dyadických škálách je cena zdvojení
$n\to 2n$ rovna $4na/(n+a)\le 4a$, tedy
\[
G_a(2^m)=4am+O_a(1),\qquad
G_a(n)=\frac{4a}{\log 2}\log n+O_a(1).
\]
Na log-log je konstantní $a$ jen svislý posun. Poměr k dolní mezi
\[
\frac{G_a(n)}{\kappa_0\log n}\to\frac{a}{a_{\mathrm{crit}}},
\qquad
a_{\mathrm{crit}}=\frac{\kappa_0\log 2}{4}\approx 0.0028479209.
\]

\begin{corollary}[nutné $a$]
Platí-li $M_n\le 1+a/n$ pro všechna dyadická $n$, pak
$a\ge a_{\mathrm{crit}}$. Existence takového $a$ není dokázána.
\end{corollary}

Na $n=2$ je $M=1$ (cena $0$) věta na $\nu_8$ a $B_*|T$.

\section{Walshův mód a Haarova hrana}

Exchangeable lift: $\widetilde P(\omega)=P_{|\omega|}/\binom4{|\omega|}$
a analogicky $\widetilde Q$.

\begin{theorem}[Walsh]
\[
\widetilde P(\omega)-\widetilde Q(\omega)=\frac{(-1)^{|\omega|}}{64}.
\tag{A}
\]
Každý marginál na nejvýše třech bitech je pro $P$ a $Q$ totožný.
\end{theorem}

\begin{proof}
Přímé dosazení po vahách $k=0,\dots,4$:
\[
\begin{aligned}
k=0&\colon& \tfrac18-\tfrac7{64}&=\tfrac1{64},\\
k=1&\colon& \tfrac1{16}-\tfrac5{64}&=-\tfrac1{64},\\
k=2&\colon& \tfrac1{24}-\tfrac5{192}&=\tfrac1{64},\\
k=3&\colon& \tfrac1{16}-\tfrac5{64}&=-\tfrac1{64},\\
k=4&\colon& \tfrac18-\tfrac7{64}&=\tfrac1{64}.
\end{aligned}
\]
Walshův charakter stupně $4$ má všechny vlastní marginály stupně $\le 3$
nulové, tedy totožné pro $P$ a $Q$.
\end{proof}

Pro tříbitový kontext váhy $r$
\[
m_r=\Bigl(\tfrac3{16},\tfrac5{48},\tfrac5{48},\tfrac3{16}\Bigr),\quad
p_r=\Bigl(\tfrac13,\tfrac25,\tfrac35,\tfrac23\Bigr),\quad
q_r=\Bigl(\tfrac5{12},\tfrac14,\tfrac34,\tfrac7{12}\Bigr).
\]

\begin{lemma}[Haarova hrana]
\[
m_r(q_r-p_r)=\frac{(-1)^r}{64},
\tag{B}
\qquad
\binom3r m_r(q_r-p_r)=\frac1{64}(1,-3,3,-1)_r.
\tag{C}
\]
\end{lemma}

\begin{proof}
Čtyři hodnoty:
\[
\begin{aligned}
r=0&\colon& \tfrac3{16}\bigl(\tfrac5{12}-\tfrac13\bigr)
&=\tfrac3{16}\cdot\tfrac1{12}=\tfrac1{64},\\
r=1&\colon& \tfrac5{48}\bigl(\tfrac14-\tfrac25\bigr)
&=\tfrac5{48}\cdot\bigl(-\tfrac3{20}\bigr)=-\tfrac1{64},\\
r=2&\colon& \tfrac5{48}\bigl(\tfrac34-\tfrac35\bigr)
&=\tfrac5{48}\cdot\tfrac3{20}=\tfrac1{64},\\
r=3&\colon& \tfrac3{16}\bigl(\tfrac7{12}-\tfrac23\bigr)
&=\tfrac3{16}\cdot\bigl(-\tfrac1{12}\bigr)=-\tfrac1{64}.
\end{aligned}
\]
Vynásobení $\binom3r$ dá (C).
\end{proof}

\section{Causal thinning na parity core}

$P=\frac78 R+\frac18 E$, $Q=\frac78 R+\frac18 O$, kde
$R=\frac1{56}(7,16,10,16,7)$,
\[
E(z)=\frac{1+6z^2+z^4}{8},\qquad
O(z)=\frac{z+z^3}{2},\qquad
E-O=\frac18(1-z)^4.
\]

\begin{theorem}[Thinning]
\[
\mathcal B_n(P,Q)\le\mathcal B_n(E,O).
\tag{D}
\]
\end{theorem}

\begin{proof}
Platí $P=\frac78 R+\frac18 E$ a $Q=\frac78 R+\frac18 O$.
V každém kroku $t$ losujeme skrytě a nezávisle $C_t\sim\mathrm{Bernoulli}(1/8)$,
nezávisle na veřejné historii.

Je-li $C_t=0$, položíme $X_t=Y_t\sim R$. Rezerva se nemění.
Je-li $C_t=1$, provedeme jeden krok exactního $(E,O)$ protokolu hloubky
nejvýše $n$: veřejný výstup má zákon $E$, vstup zákon $O$.

Pak
\[
\mathrm{Law}(Y_t\mid\mathcal F_{t-1})
=\tfrac78 R+\tfrac18 O=Q,
\]
takže vstupy jsou čerstvé $Q$. Po zprůměrování přes skryté $C_t$
\[
\mathrm{Law}(X_t\mid\mathcal H_{t-1})
=\tfrac78 R+\tfrac18 E=P,
\]
tedy veřejný marginál je $P$. Počet aktivních kroků $C_t=1$ během $n$
fyzických kroků je nejvýše $n$, proto stačí $(E,O)$ protokol hloubky $n$.
Počáteční střední rezerva $(E,O)$ stromu je přípustná i pro $(P,Q)$.
\end{proof}

Dva aktivní kroky: $K\sim\frac18(1,6,1)$, $V\sim\mathrm{Bernoulli}(1/2)$,
\[
K_1+K_2\sim\frac1{64}(1,12,38,12,1)=P_{\mathrm b},\qquad
1+V_1+V_2\sim\frac1{64}(0,16,32,16,0)=Q_{\mathrm b}.
\]
$P_{\mathrm b}-Q_{\mathrm b}=(1-z)^4/64$. To je starý coarse boundary blok.

Existuje prefix-causal vrstva
\[
\mathsf P_1\longrightarrow
\tfrac38\mathsf T+\tfrac12\boldsymbol\Phi+\tfrac18\mathscr D_2\mathsf P_1.
\tag{E}
\]
Pro $\mathcal C=\{b_0=0,\ b_1\le 13/16\}$:
\[
\tfrac12\longrightarrow\tfrac{15}{38}\longrightarrow\tfrac{189}{722}
<\tfrac{13}{48}\longrightarrow\mathcal C,
\]
nejvýše tři critical kroky, repair margin alespoň $7/128$.

\section{Blokový most}

$E^{*2}(z)=P_{\mathrm b}(z^2)$, $O^{*2}(z)=Q_{\mathrm b}(z^2)$.
Seskupením dvojic symbolů ve stromu hloubky $2n$ pro $(E,O)$ vznikne
strom hloubky $n$ pro $(E^{*2},O^{*2})$; dělení rezervy dvěma ho
převádí na coarse pair.

\begin{theorem}[blokový most]
\[
2\,\mathcal B_n(P_{\mathrm b},Q_{\mathrm b})
\le
\mathcal B_{2n}(E,O)
\le
2\,\mathcal B_n(P_{\mathrm b},Q_{\mathrm b})+4.
\tag{*}
\]
S thinningem (D)
\[
\mathcal B_{2n}(P,Q)
\le
2\,\mathcal B_n(P_{\mathrm b},Q_{\mathrm b})+4.
\tag{**}
\]
\end{theorem}

\begin{proof}
Generující funkce dávají $E(z)^2=P_{\mathrm b}(z^2)$ a $O(z)^2=Q_{\mathrm b}(z^2)$:
\[
(1+6z^2+z^4)^2=1+12z^2+38z^4+12z^6+z^8,
\]
což je $64\,P_{\mathrm b}(z^2)$, a analogicky pro $O$.

Nechť $T$ je exactní $(E,O)$ strom hloubky $2n$ s kořenovou rezervou $J$.
Seskupme po sobě jdoucí dvojice symbolů. Součet dvou $E$-výstupů leží
na sudé mřížce $\{0,2,4,6,8\}$ a po dělení dvěma má zákon $P_{\mathrm b}$.
Součet dvou $O$-vstupů po dělení dvěma má zákon $Q_{\mathrm b}$.
Rozdíl rezervy přes dvojici je sudý. Položíme-li $J'=J/2$ na
půlškále, dostaneme exactní $(P_{\mathrm b},Q_{\mathrm b})$ strom hloubky
$n$. Proto $\E J\ge 2\mathcal B_n(P_{\mathrm b},Q_{\mathrm b})$, což je
levá nerovnost v (*).

Opačně: z exactního $(P_{\mathrm b},Q_{\mathrm b})$ stromu hloubky $n$
zvedneme rezervu faktorem $2$ a každý coarse krok nahradíme dvojicí
$(E,O)$ kroků. Zpoždění nejvýše jednoho coarse kvanta velikosti $4$
(maximální výstup páru po půlení) dává pravou nerovnost.
Nerovnost (**) je (*) plus (D).
\end{proof}

Tedy
\[
\mathcal B_n(P,Q)=\Theta(\log n)
\quad\Longleftarrow\quad
\mathcal B_n(P_{\mathrm b},Q_{\mathrm b})=O(\log n).
\]
Ordered-history lift z coarse bloku na symboly už není otevřená mezera.

Zkratka $B_1=B_3=B$ dává
$P_c=\frac1{40}(1,0,38,0,1)$, $Q_c=\frac1{40}(0,4,32,4,0)$,
$P_c=\frac45\delta_2+\frac15 E$, $Q_c=\frac45\delta_2+\frac15 O$.
Hmotnost $1/5$ není horizon $n/5$. Není rekurzivní důkaz.

\section{Weak $L^1$ --- převod, ne věta}

Předpokládejme regeneraci $\Phi$ s
\[
m_s\le C\,2^{-s},\qquad J\le K\,2^s.
\tag{F,G}
\]
Pak $\E\min(J,N)\le (CK/\log 2)\log N+O(1)$. Antecedent není dokázán.
To je algebra ``pádu atmosférou''.

Nosný otevřený objekt:
\[
\Phi_s\stackrel{\mathrm{causal}}{\longrightarrow}
\text{finite boundary grammar}+\tfrac12\,\Phi_{s+1}.
\]

\section{Dyadický boundary-flow}

Pětivětvový fan $(P_{\mathrm b},Q_{\mathrm b})$. Na critical core
$\mathcal C=\{b_0=0,\ b_1\le 13/16\}$ platí $B_0,B_2\in\mathcal C$ a hard edge
\[
\tfrac12\to\tfrac{15}{38}\to\tfrac{189}{722}
<\tfrac{13}{48}\to\mathcal C,
\]
s repair marginem alespoň $7/128$.

Faktor $1/2$ \textbf{není} vlastnost $n\to n-1$. Půlení na každém kroku
stromu by dalo $2^{-t}$, ne $1/n$. Účet
\[
p_n=\frac1{8n},\qquad p_{2n}=\frac12 p_n,
\qquad
\frac{1/(16n)}{1/(8n)}=\frac12
\]
říká, že $1/2$ je survival přes slupku $4n\to 8n$. Pět dětí packet jen nese
kauzálně. Zakázáno je
\[
\forall x:\quad \mathrm{top}(B_x)\le\tfrac12\,\mathrm{top}(\beta_{n-1}).
\]

\begin{proposition}[prefixový zámek]
Fan v uzlu $h$ se volí před všemi dětmi $hx$. Nelze dát pokračování
$0^{t+1}$ jiný donor než $0^t2$. To je podmínka exactnosti (ZR), ne
konstrukce.
\end{proposition}

\begin{conjecture}[dyadic boundary-flow]
Existuje konstanta $C<\infty$ a jeden prefix-kompatibilní, phase-indexed
fan se stejnými donor packets pro všechny děti naráz, který realizuje
\[
\mathcal L_n(\delta)\longrightarrow
\mathcal L_{2n}\!\left(\tfrac12\delta+C\delta^2\right).
\]
Ekvivalentně: boundary packet $p_n\delta_{4n}$ lze přes depth-$n$ causal
fan transportovat do dyadické slupky tak, že každý leaf je znovu
admissible na škále $2n$ a pouze podíl $\frac12+O(\delta)$ zůstává
kritický.
\end{conjecture}

Na stěně $b_0=0$ je hlášen exactní invariant
$W(B)=R_5(B)-4b_1$, $W(C_0)=W(B)+\frac18 b_4$. Aktivní větve $x=2,4$ při
pullbacku vytáhnou $b_6$ a další koeficienty. Pozitivní konečná ani
countable gramatika se z nich neuzavřela. Tento článek \textbf{nedokazuje}
Conjecture výše.
\[
\boxed{\text{NEUZAVŘENO}}
\]
To není protiargument proti $\mathcal B_n=O(\log n)$. Chybí pozitivní
substitution lemma.

\section{Další dokázané kotvy z ledgeru}

\begin{itemize}
\item $H(B)\le 8/11$ pro nekonečně pokračovatelné stavy,
$H(B)=B(0)+(7/27)B(1)$.
\item $B\in K_2\iff H(B)\le 200/243$.
\item $b_0\le 13/20\Rightarrow B\succeq_{\mathrm{st}}\rho_4\Rightarrow B\in K_4$.
\item $\Xi_{2,4}(B_h;1)=0$ na $\nu_8$ a $B_*|T$. Tedy $n=2$, $M=1$.
\item $\nu_8\in K_8$, $\beta_4\in K_8$, $\beta_8\in K_8$.
\item $\Gamma_5(\mu)=531441/446731$. Pevné semínko $\mu$ nemá
$\Gamma_n=1+O(1/n)$.
\item Afinní $\Xi_{1,s}(\mu;M)$ na dokumentovaných intervalech $M$.
\end{itemize}

\section{Zakázáno}

\begin{itemize}
\item $Q*\beta_n\not\succeq_{\mathrm{st}}P*\beta_{n-1}$ (gap $5/256$ v $k=1$).
\item $\{p\le 1/4\}$ není invariant: $p=1/4$, $q=3/4$ dává $p'_2=11/32$.
\item $\{B\succeq B^{(3)}\}$ není invariant: $T_2(B^{(3)})\not\succeq_{\mathrm{st}}B^{(3)}$.
\item $\{h_{r+1}\le h_r\}$ a $\{h_1\le q\}$ padají na $T_0^{18}(B^{(3)})$.
\item Stacionární interior klec HEAD-max na ose $b_1=0$ neexistuje.
\end{itemize}

\section{Otevřené}

Jediný theorem-level krok k $\mathcal B_n=O(\log n)$ je Conjecture
(dyadic boundary-flow): jeden prefix-kompatibilní fan. Z mostu (*) to
stačí. Dále zůstávají (Q) pro $\beta_{2^m}$, barrier front po $T_2/T_4$
a existence $B_*\in K_\infty$. Žádné z nich tento článek neprohlašuje
za větu.

\section{Závěr}

Dokázáno je
\[
\kappa_0\log n-o(\log n)
\le\mathcal B_n
\le\frac{473}{56}\,n^{1/4},
\]
Walshův mód, thinning (D) a blokový most (*).
\textbf{Neštěpíme hmotu mezi větve. Rozbalujeme ji mezi škály.}
Ten unwrap není sestrojen. Proto
\[
\boxed{\mathcal B_n=O(\log n)\quad\text{je otevřené.}}
\]

\end{document}
